% ==============================================================================
% =================================== Aufgabe 2 =================================
% ==============================================================================

\chapter{Aufgabe}
\label{chap:2 Aufgabe}
\thispagestyle{fancy}
\section{Themeneingrenzung \& Hypothese}
\begin{enumerate}[label=\thechapter.\arabic*), leftmargin=*]
    \item \textbf{Themeneingrenzung} \par Das Oberthema „\textbf{\textit{Rollen im Vergleich von klassischem und agilem Projektmanagement}}“ wird auf folgende Forschungsfrage eingegrenzt: \par „\textbf{\textit{Der Einfluss des Rollenwechsels vom hierarchischen Projektleiter zum agilen Scrum Master auf die intrinsische Motivation von Softwareentwicklungsteams}}“.
    \item \textbf{Begründung der Eingrenzung} \par Eine vollständige Gegenüberstellung aller Rollen (\textit{Auftraggeber}, \textit{Team}, \textit{Product Owner \acs{etc.}}) würde den geforderten Umfang von 40–60 Seiten sprengen und zu oberflächlich bleiben. Die Eingrenzung auf die Führungsrolle ist sinnvoll, da Dipl.-Pädagoge Bernd-Uwe Kiefer \cite[1]{1} explizit betont, dass die „\textbf{\textit{Soft Facts}}“ (\textit{wie Führung und Motivation}) über den Projekterfolg entscheiden, während „\textbf{\textit{Hard Facts}}“ (\textit{Strukturen}) diesen lediglich unterstützen. Der Fokus auf die \acs{IT}-Branche legitimiert sich dadurch, dass agile Methoden wie Scrum dort ihren Ursprung haben und am weitesten verbreitet sind.
    \item \textbf{Zielsetzung und Hypothese} 
    \begin{itemize}
        \item \textbf{Zielsetzung:} Das Ziel der Arbeit ist es, die Veränderung der Motivationsstrukturen zu analysieren, wenn die zentrale Steuerungsgewalt eines klassischen Projektleiters durch das Prinzip des „\textbf{\textit{Servant Leadership}}“ eines Scrum Masters ersetzt wird.
        \item \textbf{Hypothese:} Der Übergang von der zentralen Steuerung durch einen klassischen Projektleiter hin zur unterstützenden Moderation durch einen Scrum Master korreliert positiv mit der intrinsischen Motivation des Projektteams.
    \end{itemize}
    \item \textbf{Begründung anhand wissenschaftlicher Qualitätskriterien} \par Die gewählte Zielsetzung und Hypothese erfüllen folgende Kriterien:
    \begin{enumerate}[label=\thechapter.\arabic{enumi}.\arabic*), leftmargin=*]
        \item \textbf{Relevanz:} Das Thema ist in den aktuellen Megatrend der „\textbf{\textit{agilen Transformation}}“ eingebettet und besitzt somit eine hohe wissenschaftliche und praktische Bedeutung.
        \item \textbf{Validität (\textit{Gültigkeit}):} Durch die Konzentration auf zwei skalierbare Variablen (\textit{Führungsstil als Variable A und Motivation als Variable B}) wird sichergestellt, dass die Arbeit tatsächlich den beabsichtigten Zusammenhang untersucht.
        \item \textbf{Reliabilität (\textit{Zuverlässigkeit}):} Die Untersuchung basiert auf klar definierten Rollenprofilen nach IMI03 (\textit{Projektleiter \acs{vs.} Scrum Master}), wodurch die Ergebnisse bei einer Wiederholung unter gleichen Bedingungen reproduzierbar sind.
    \end{enumerate}
\end{enumerate}

% ==============================================================================
% =================================== Aufgabe 2 =================================
% ==============================================================================